% Add "answers" option to print the answers
\documentclass[addpoints]{exam}

\usepackage[utf8]{inputenc}
\usepackage{calc}
\usepackage[dvipsnames]{xcolor}
\usepackage{ifthen}
\usepackage{tcolorbox}

% \makebox[\textwidth]{Name:\enspace\hrulefill}
%
% \vspace{5mm}
%\makebox[\textwidth]{Instructor’s name:\enspace\hrulefill}

\definecolor{MyLightGray}{gray}{0.60}

\newcommand{\truefalse}
{		\begin{checkboxes}
			 \choice True
			 \choice False. Explanation:
		\end{checkboxes}
		\vspace{10pt}
}


% Check if 'answers' is an option
\newcommand{\ifanswers}[2]{%
  \ifthenelse{\boolean{exam@answers}}{#1}{#2}%
}



\newcommand{\answerspace}[1]
{
	\ifprintanswers
	\else
		\setlength{\mylength}{10pt}%
		% \vspace{#1\mylength}
		% Make the grey box
		\par\nobreak\noindent
		\makebox[\linewidth][l]{%
			% set inner padding to zero, so only the frame shows
			{\setlength{\fboxsep}{0pt}%
			% draw a color frame, white (or transparent) inside
			\fcolorbox{MyLightGray}{white}{%
			% the inner box has zero width content but fixed height
			\parbox[t][#1\mylength][c]{%
				% subtract twice the frame‐thickness from the full width
				\dimexpr\linewidth-2\fboxrule\relax
			}{\strut}%
			}%
			}%
		}%
		\par
		\vspace{3\mylength}

  	\fi
}
\newlength{\mylength}

\newcommand{\code}[1]{{\texttt{#1}}}


%\usepackage{listings}
\usepackage{color}
\usepackage{lastpage}
\usepackage{verbatim}
\usepackage{fancyvrb}
\usepackage{graphicx}
\usepackage{float}
\usepackage{minted}
\usepackage{multicol}
\usepackage{enumitem}

\newcommand{\includecode}[2][c]{\lstinputlisting[caption=, escapechar=, language=#1]{#2}}
\RecustomVerbatimCommand{\VerbatimInput}{VerbatimInput}{}

\newcommand{\Code}[1]{\mintinline{python}{#1}}

\renewcommand{\solutiontitle}{\noindent\textbf{Solution sketch: }}

\definecolor{SolutionColor}{rgb}{0.8,0.9,1}

\cfoot{Page \thepage\ of \pageref{LastPage}}

\rhead{Initials: \rule{20mm}{0.4pt}}


%%%%%%%%%%%%%%%%%%%%%%%%%%%%%%%%%%%%%%%%
\begin{document}
\fvset{numbers=left,numbersep=3pt}
\thispagestyle{empty}

\begin{center}
	{\LARGE CS 408 Fall 2026 -- Input Generators Quiz \\}
\fbox{\fbox{\parbox{5.5in}{\centering
	\begin{itemize}
        \item Try to come up with a solution on your own first, but
          after that talk to your friends and neighbors.

        \item Full points will be rewarded for \textbf{any} (reasonable)
          attempt.
        \end{itemize}
	}}}
\end{center}

\makebox[\textwidth]{\textbf{Name (please print)}: \rule{100mm}{0.8pt} }

\vspace{10mm}
\makebox[\textwidth]{\textbf{PUID}: \rule{100mm}{0.8pt} \hspace{35mm}}

\begin{questions}

\question We want to test an implementation of binary search trees, so we
  need a supply of them. Represent a tree as either \Code{None}, for an
  empty tree, or a triple \Code{(left, value, right)}. Recall that in a
  \textit{binary search tree} (BST) every value in the left subtree is smaller
  than the node's value, and every value in the right subtree is larger.

  Here is a generator that produces binary trees, but does not
  guarantee they are valid BSTs:

\begin{minted}{python}
@st.composite
def trees(draw, lo=-1000, hi=1000, max_depth=5):
    if max_depth <= 0 or draw(st.integers(0, 2)) == 0:
        return None
    return (draw(trees(lo, hi, max_depth - 1)),
            draw(st.integers(min_value=lo, max_value=hi)),
            draw(trees(lo, hi, max_depth - 1)))
\end{minted}

\begin{parts}

\part Why would it be a bad idea to generate trees using \Code{trees}
  and then filter out any that are not valid BSTs? %
  \answerspace{16} %
  \begin{solution} %
    Almost none of them are. Sampling 3000 trees from a generator like
    this one turns up search trees with three or more nodes about
    \textbf{2\%} of the time, and the rate collapses as the trees grow
    --- a tree of $n$ distinct values has $n!$ possible arrangements and
    only a handful of them are ordered correctly.

    So filtering wastes almost all of the work, and worse, it
    \emph{biases} what survives: the trees that get through are
    overwhelmingly the tiny ones, because those are the only ones with a
    decent chance of being ordered by accident. The tests would pass
    while hardly ever exercising a tree of interesting size.

    (This is the same trap as \Code{assume(is_sorted(xs))} from earlier:
    the health check fires, or --- worse --- it does not, and you are
    quietly testing almost nothing.)
  \end{solution}

\newpage

\part Modify \Code{trees} so that it produces \emph{only} valid
  BSTs. You do not need to get the Hypothesis syntax exactly right,
  but try to be precise about what each recursive call produces, and
  why that guarantees the correctness of the modified generator. %
  \answerspace{30} %
  \begin{solution} %
    Draw the node's value first, then \textbf{narrow the bounds} handed
    to each subtree so that the ordering invariant cannot be broken:

\begin{minted}{python}
@st.composite
def bsts(draw, lo=-1000, hi=1000, max_depth=5):
    if max_depth <= 0 or lo > hi or draw(st.integers(0, 2)) == 0:
        return None
    v = draw(st.integers(min_value=lo, max_value=hi))
    return (draw(bsts(lo, v - 1, max_depth - 1)),
            v,
            draw(bsts(v + 1, hi, max_depth - 1)))
\end{minted}

    Everything in the left subtree is drawn from $[lo, v-1]$ and
    everything in the right from $[v+1, hi]$, so no comparison can come
    out wrong. Nothing is ever rejected.

    The \Code{lo > hi} test matters: once the interval is empty there is
    no value left to draw, so the subtree must be \Code{None}.

    Another correct answer: generate a sorted list of distinct values and
    build a tree from it. That works, but think about which shapes of
    tree it can produce --- if it always splits at the midpoint, it will
    only ever hand you balanced trees.
  \end{solution}

\end{parts}

\end{questions}

\end{document}
